\documentclass[12pt]{article}
\usepackage{amsmath,amssymb}
\usepackage[utf8]{inputenc}
\usepackage[T2A]{fontenc}
\usepackage[russian]{babel}
\usepackage{enumitem}
\usepackage[a4paper,margin=2cm]{geometry}
\usepackage{titlesec}

\titleformat{\section}[block]{\large\bfseries\centering}{}{1em}{}
\titleformat{\subsection}{\normalfont\bfseries}{\thesubsection.}{0.5em}{}

\title{XXX Всероссийская олимпиада по астрономии\\
Заключительный этап, ФТ «Сириус», 27 марта -- 2 апреля 2023 г.\\
Теоретический тур}
\date{}

\begin{document}

\maketitle

\section*{9 класс}

\subsection{Свидание с Венерой}
Ближайшее нижнее соединение Венеры с Солнцем по эклиптической долготе произойдёт
13 августа 2023 года в $11^h10^m$ по Всемирному времени. Известно, что координаты
Солнца в этот момент составят $\alpha = 9^h31.5^m$, $\delta = +14^\circ40'$, а Венера
пройдёт в $7^\circ40'$ южнее эклиптики. Определите координаты точки на поверхности
Земли, из которой Венера будет лучше всего видна в этот момент. Считать, что Венера
видна, если центр диска Солнца расположен не выше горизонта, а критерием качества
видимости при этих условиях является высота Венеры над горизонтом. Атмосферной
рефракцией и уравнением времени пренебречь.

\subsection{Семейство Хильды}
Около орбиты Юпитера находится группа астероидов, образующих так называемое
семейство Хильды. Эти астероиды примечательны тем, что в любой момент времени
образуют правильный треугольник, вершины которого лежат вблизи орбиты Юпитера. Этот
треугольник поворачивается в пространстве синхронно с орбитальным движением Юпитера
так, что Юпитер всегда равноудалён от двух ближайших к нему вершин. Определите
большую полуось (с точностью лучше 1\%) и эксцентриситет орбиты астероидов из
семейства Хильды. Орбиту Юпитера считайте круговой, «толщиной линий» треугольника
пренебрегите.

\subsection{Триангуляция будущего}
В ходе астрометрической миссии будущего (через 200 с небольшим лет) координаты
центра сверхмассивной чёрной дыры (СМЧД) в центре нашей Галактики одновременно
измеряются инфракрасными телескопами с околоземной орбиты и с поверхности Плутона
(его координаты в небе Земли в этот момент $\alpha = 18^h$, $\delta = -17^\circ$,
гелиоцентрическое расстояние 32~а.\,е.). Обоими телескопами было сделано по одному
измерению с точностью до 1 наносекунды дуги ($10^{-9}{}''$). С какой точностью можно
определить расстояние до центра СМЧД в центре Галактики на основе этих измерений?

\subsection{Улетающая звезда}
В 2014 году была открыта звезда WISE 0720$-$0846. Эта звезда интересна тем, что в
прошлом она достаточно близко подлетела к Солнечной системе. Определите:
\begin{enumerate}
  \item Минимальное гелиоцентрическое расстояние, на которое сближалась звезда с Солнцем.
  \item Сколько лет назад это произошло?
  \item Какова была при этом её звёздная величина?
  \item Чему были равны её полное собственное движение и радиальная скорость в этот момент?
  \item Если полагать, что данная звезда породила возмущение в кометном облаке Оорта,
        то через какой минимальный промежуток времени следует ожидать приток комет в
        окрестностях Земли?
\end{enumerate}

Современные характеристики звезды:
\begin{center}
\begin{tabular}{lcc}
Видимая звёздная величина & $m$ & 18.3 \\
Лучевая скорость & $v_R$ & $+82.4$~км/с \\
Параллакс & $\pi$ & $0.147''$ \\
Собственное движение вдоль экватора & $\mu_\alpha$ & $-0.0403''/\text{год}$ \\
Собственное движение к полюсу & $\mu_\delta$ & $+0.1148''/\text{год}$ \\
\end{tabular}
\end{center}

\subsection{Горячее будущее}
Согласно модели, используемой в статье K.P.~Schröder, R.C.~Smith (2008), через
12.17 млрд лет после своего образования Солнце достигнет наибольшего размера в
течение своей эволюции. Радиус Солнца будет в 256 раз больше нынешнего, светимость
станет в 2730 раз больше нынешней, а масса уменьшится на 33.2\%. Считайте, что потеря
массы Солнцем происходит медленно в течение всей стадии красного гиганта, а
физические свойства других тел Солнечной системы при росте температуры не меняются.

Определите для момента времени, описанного выше:
\begin{enumerate}
  \item Какие планеты Солнечной системы будут поглощены Солнцем?
  \item Найдите среднюю температуру поверхности самой горячей планеты, которая не
        будет поглощена Солнцем.
  \item Какие крупные тела Солнечной системы окажутся в зоне жизни (средняя
        температура от 250 до 300~K без учёта парникового эффекта в атмосфере)?
\end{enumerate}

\subsection{Далёкая галактика}
Диск далёкой спиральной галактики расположен «плашмя» по отношению к лучу зрения.
Характерная величина поверхностной яркости диска спиральной галактики $20^m$ с
квадратной секунды. Определите характерное количество звёзд в 1~пк$^3$ диска.
Межзвёздным поглощением света пренебречь.

\clearpage
\section*{10 класс}

\subsection{Пролетая над столицами}
Спутник, движущийся по круговой орбите вокруг Земли, последовательно побывал в
зените в Каракасе ($10^\circ30'$~с.ш., $66^\circ56'$~з.д.) и Кито ($0^\circ13'$~ю.ш.,
$78^\circ31'$~з.д.) с интервалом 20 минут. Определите радиус и наклонение орбиты
спутника.

\subsection{В закатном небе}
Во время недолгого перерыва в работе строители новой научной базы на поверхности
Марса любуются закатом Солнца. Сразу после него на фоне зари появились две яркие
планеты -- Венера и Земля, вступившие в тесное соединение (без покрытия) друг с
другом невдалеке от Солнца на небе и имеющие одинаковую видимую яркость. Определите
угловое расстояние Венеры и Земли от Солнца, а также величины их фаз. Считать орбиты
всех планет круговыми. Принять также, что за счёт плотных атмосфер и облаков
поверхностная яркость освещённых частей дисков Венеры и Земли не зависит от фазового
угла (между направлением на Солнце и наблюдателя). Сумеречными эффектами
(освещённостью атмосфер Венеры и Земли во время сумерек) пренебречь.

\subsection{Кубодетектор}
В жёсткой гамма-области, где невозможно построение телескопа как оптической схемы,
направление на источник иногда определяется по отношению потоков энергии от него
через площадки детектора, по-разному ориентированные в пространстве. В космос
запускается гамма-детектор, имеющий форму куба, каждая из шести граней которого
имеет одинаковую чувствительность и фиксирует поток энергии через свою площадь с
внешней стороны куба. Она это делает с точностью 3\%, то есть при истинном потоке
$J$ измеренный с равной вероятностью будет лежать в интервале от $0.97J$ до $1.03J$.
Определите максимально возможную угловую «площадь ошибок» (в квадратных градусах)
участка неба, где может находиться источник, координаты которого измерены детектором
один раз.

\subsection{Таинственный мир}
Экспедиция прибыла на планету, обращающуюся вокруг далёкой звезды, и приступила к
исследованиям. Оказалось, что в месте базирования экспедиции местная звезда каждый
день проходила через зенит, но климат существенно менялся в течение года. Участники
экспедиции заметили, что ночью комфортные условия ($T=+20^\circ$C) в палатке без
обогрева в тёплый сезон достигались при длительном нахождении там одного человека, а
в холодный сезон для этого требовалось присутствие сразу трёх человек. При этом
«солнечные» сутки в тёплый сезон были на 0.2\% длиннее, чем в холодный, а местный год
состоял из 100 местных «солнечных» суток. Перепад температур между ночью и днём был
постоянен в течение года и составлял $20^\circ$. Атмосфера была вполне пригодной для
дыхания, при этом была очень сухой и не обладала парниковыми свойствами, поверхность
состояла в основном из очень тёмных пород, имеющих малую теплопроводность.
Спектральный состав излучения звезды и её физические свойства были аналогичны
солнечным. Определите период осевого вращения планеты и эксцентриситет её орбиты.

\subsection{Урановое Солнце}
На заре развития ядерной физики было высказано предположение, что энерговыделение
Солнца в течение многих миллиардов лет обеспечивалось распадом урана-238. Известно,
что период полураспада этого изотопа урана составляет 4.47 млрд лет, что можно
считать равным возрасту Земли, который был уже известен к тому времени. При этом
образуются другие радиоактивные изотопы, распад которых происходит существенно
быстрее. Энерговыделение в ходе всего цикла реакций, начиная с одного атома
$^{238}$U, составляет 6.7~МэВ. Определите минимально возможную массу Солнца в
предположении такого механизма его свечения.

\subsection{Гигантский обзор неба}
В 2024 году планируется к запуску проект LSST (Large Synoptic Survey Telescope), или
телескоп имени Веры Рубин, построенный в Чили (широта $-30^\circ$). Одна из основных
задач этого телескопа -- автоматизированный поиск сверхновых звёзд. Телескоп имеет
составное зеркало эффективным диаметром 6.7 метра и поле зрения 9.6 квадратного
градуса. Телескоп оснащён системой адаптивной оптики, исправляющей атмосферные
искажения, и системой ПЗС-матриц общим объёмом 3.2 гигапикселя. Авторы проекта
считают, что телескоп может делать полный обзор всей видимой из обсерватории части
неба (высота объекта над горизонтом не менее $55^\circ$) за 3 ночи и за 10 лет
открыть 3 миллиона сверхновых звёзд. На основании этих данных оцените количество
спиральных галактик на кубический мегапарсек.

Считайте, что сверхновые звёзды вспыхивают в основном в спиральных галактиках, и это
происходит в каждой из них в среднем раз в 100 лет. Фон неба в месте постройки
обсерватории составляет $21.5^m$ на квадратную секунду, сверхновую можно обнаружить,
если средние отсчёты от неё хотя бы на одном пикселе матрицы превосходят отсчёты от
фона. Межзвёздным поглощением и влиянием космологических факторов на видимую яркость
далёких объектов пренебречь.

\clearpage
\section*{11 класс}

\subsection{«В далёком созвездии Тау Кита»}
В планетной системе около звезды $\tau$~Кита с массой 0.78 массы Солнца обращается
несколько планет. Планета $c$ обладает круговой орбитой радиуса 0.195~а.\,е.,
планета $e$ движется в той же плоскости и в том же направлении по эллиптической
орбите с большой полуосью 0.538~а.\,е. и эксцентриситетом 0.18. В некий момент
времени произошло великое противостояние планеты $e$ для наблюдателя на планете $c$.
Какой будет фаза планеты $c$ при наблюдении с планеты $e$ и фаза планеты $e$ при
наблюдении с планеты $c$ спустя 10 дней после великого противостояния?

\subsection{Метеорное эхо}
Радар, изучающий метеоры, фиксирует радиоэхо только тогда, когда направление от
радара к метеору перпендикулярно ионизационному следу метеора. Радар, расположенный
на широте $+52^\circ$, направлен в точку востока. В $11^h02^m$ местного времени
2 июля он начинает получать многочисленные сигналы от метеоров, проходящих через его
луч. Радар поворачивают на $40^\circ$ вдоль горизонта в сторону юга. В новом
положении он начинает обнаруживать метеоры того же потока в $12^h42^m$. Найдите
экваториальные координаты (прямое восхождение и склонение) радианта этого метеорного
потока. Уравнением времени пренебречь.

\subsection{Инфракрасная камера}
ИК-камера ASTRONIRCAM 2.5-м телескопа Кисловодской горной обсерватории ГАИШ МГУ не
позволяет наблюдать в фильтре $J$ звёзды ярче $9.0^m$ из-за того, что при таком
блеске поток фотонов за минимально возможное для камеры время накопления полностью
заполняет ячейки детектора, и дальнейшая регистрация фотонов становится невозможной.
Оцените аналогичную предельную звёздную величину для фильтра $K$. Детектор камеры
обладает практически постоянным квантовым выходом во всём рабочем интервале длин
волн, охватывающем обе полосы. Считать, что фильтр $J$ полностью пропускает излучение
с длинами волн от 1.17 до 1.34~мкм, а фильтр $K$ -- от 2.04 до 2.35~мкм, не пропуская
излучение вне этих интервалов. Учесть, что звёздная величина звезды Вега (температура
10000~K) равна $0^m$ в обеих спектральных полосах.

\subsection{Горячая пыль}
Измерения в ИК-диапазоне показали, что белый карлик окружает кольцо из тёмных
пылинок с температурой не более 1500~K. Предполагается, что пылинки попадают в
кольцо в результате приливного разрушения пролетающих мимо астероидов. Оцените
внутренний и внешний радиусы кольца в километрах, если масса белого карлика 0.8 масс
Солнца, радиус белого карлика 0.01 радиусов Солнца, его температура 20000~K, а
плотность астероидов 1~г/см$^3$.

\subsection{Неразумное бегство}
Космический аппарат массой $m$ обращается по круговой околосолнечной орбите с
радиусом $R_0$. В один момент времени аппарат включает двигатель, создающий
постоянную силу тяги, всё время направленную от Солнца. При какой минимальной
величине этой силы аппарат в итоге покинет Солнечную систему? Считать, что масса
аппарата в результате работы двигателя не меняется, взаимодействием с другими телами
Солнечной системы, кроме Солнца, и излучением пренебречь.

\subsection{Волчья Сверхновая}
В мае 1006 года наблюдатели на Земле зафиксировали сильнейшую вспышку звезды в точке
неба, находящейся в современном созвездии Волка. По всей вероятности, это была
Сверхновая типа Ia, самая яркая на Земле с начала новой эры. Её видимая звёздная
величина в максимуме оценивается в $-7.5^m$. Расстояние до неё составляло 2.2~кпк.

Как известно, кривая блеска сверхновых звёзд типа Ia после максимума определяется
процессом распада никеля-56, образующегося при вспышке:
\[
{}^{56}\mathrm{Ni} + e^- \rightarrow {}^{56}\mathrm{Co} + 1.75~\text{МэВ}
\quad (\text{период полураспада } 6.10 \text{ суток});
\]
\[
{}^{56}\mathrm{Co} + e^- \rightarrow {}^{56}\mathrm{Fe} + 3.61~\text{МэВ}
\quad (\text{период полураспада } 77.12 \text{ суток}).
\]

Исходя из этих данных, определите общую массу никеля-56, образовавшегося при вспышке
сверхновой, и время после регистрации максимума, через которое звезда в небе Земли
ослабла до $0^m$. Сверхновая наблюдалась в стороне от Млечного пути, межзвёздным
поглощением света и болометрической поправкой сверхновой пренебречь.

\end{document}
